\documentclass[10pt,a4paper]{amsart}
\usepackage[margin=19mm]{geometry}
\usepackage{amsmath,amssymb,mathtools}
\usepackage[scr=rsfs,cal=euler]{mathalfa}
\usepackage{xfrac}
\usepackage[unicode,hidelinks,bookmarks=false]{hyperref}
\newcommand{\ie}{i.e.\ }
\newcommand{\cf}{cf.\ }
\newcommand{\resp}{respectively\ }
\newcommand{\Subsection}{\subsection}
\providecommand{\nobreakdash}{\nobreak\hbox{-}}
\setlength{\emergencystretch}{2em}
\multlinegap0pt
\usepackage{amssymb}
\newtheorem{lemma}{Lemma}
\newcommand{\C}{\mbb{C}}
\newcommand{\ann}{\mfrak{A}}
\renewcommand{\bar}{\overline}
\newcommand{\basgen}{\mcal{B}}
\newcommand{\class}[1]{\left[#1\right]}
\newcommand{\intt}{\mfrak{I}}
\newcommand{\lift}[1]{\tilde{#1}}
\newcommand{\lp}{\left(}\newcommand{\rp}{\right)}
\newcommand{\mbb}[1]{\mathbb{#1}}
\newcommand{\mcal}[1]{\mathcal{#1}}
\newcommand{\mfrak}[1]{\mathfrak{#1}}
\newcommand{\mrm}[1]{\mathrm{#1}}
\newcommand{\qroot}{q^{1/n}}
\newcommand{\rev}[1]{\bar{#1}}
\newcommand{\sC}{\C^\ast}
\newcommand{\sk}{\mcal{S}_n^q}
\newcommand{\sub}[1]{\mrm{sub}\lp#1\rp}
\newcommand{\sym}{\mrm{Sym}}
\renewcommand{\tilde}{\widetilde}
\title{Monomial web basis for the SL(N) skein algebra of the twice punctured sphere}
\author{Tommaso Cremaschi, Daniel C. Douglas}
\date{Source: arXiv:2407.04178v3 (CC BY 4.0)}
\begin{document}
\maketitle

\noindent\emph{Source and licence.} Monomial web basis for the SL(N) skein algebra of the twice punctured sphere. Version arXiv:2407.04178v3 (CC BY 4.0), distributed by arXiv under the Creative Commons Attribution 4.0 International licence, http://creativecommons.org/licenses/by/4.0/, as recorded in the arXiv licence metadata for this exact version and shown on the abstract page https://arxiv.org/abs/2407.04178v3. Full licence notice: \texttt{licenses/arxiv-2407.04178-CC-BY-4.0.txt}.\par\smallskip

\noindent\textit{Editorial scope.} Counted unit: the article's lemma lem:inductionlargei (source0.tex lines 527--529). The article's complete original proof of it is delivered with it (source0.tex lines 531--566, one \texttt{proof} environment of 36 lines). No mathematics is written, repaired or re-derived: the statement and the proof are the article's own lines, in the article's own notation, and no retained line is substituted or re-flowed.  Retained uncounted as the source-local context the proof needs: lines 434--439; lines 443--447; lines 514--515; lines 520--525.  One declared adaptation: exactly 5 ancillary strings inside the retained spans are omitted (image-inclusion commands and, where the source repeats a label, the repeated label definitions) and no other line differs from the source; the figure environments, with their placement, captions and labels, are kept, and the package records the omitted strings in fidelity receipts bound to the affected bodies.  No retained line contains a citation, so the wrapper carries no bibliography and no bibliography entry.  The retained lines contain the article's own diagram code, which the wrapper typesets with the standard tikz packages; no image file is delivered and the package declares no asset dependency.  Editorial formatting only, in addition: an AMS \texttt{amsart} preamble that restates the article's own declarations and macros used by the retained lines, namely the environments \texttt{lemma} on separate counters, together with the standard packages those lines need.  The retained environments print in this document as Lemma 1 in order of appearance, whereas the article prints the same items with the numbering of its own full document.  The complete native source of the article as distributed in the arXiv e-print for this exact version is bundled unchanged as \texttt{sources/161839/source0.tex}.
\begin{figure}[htb!]
\centering

\caption{Algebraic intersection number.}
\label{fig:fig7}
\end{figure}

\begin{lemma}\label{lem:refinement}
Let $L$ be a link in $\ann\times\intt$ and let $\gamma$ be the projection of $L$ on $\ann$.  If $\gamma$ has at most $i$ positive intersections and $j$ negative intersections with an arc $\alpha$ oriented out of the puncture and cutting the annulus, colored blue 
%depicted as a dotted ray
in Figure {\upshape\ref{fig:fig7}}, then $\class{L}\in\sub{\gamma_m;-j\leq m\leq i}$.  
\end{lemma}

\begin{lemma}\label{lem:square-lemma}
Let $a=q^{n(n-1)}$, $b=-q^{(1-n)/n}$, and $a_k=q^{-k}[n-k]$.  For all $\qroot\in\sC$, the relation appearing in Figure {\upshape\ref{fig:square-lemma}} holds.  The same relation holds with all orientations reversed.\end{lemma}

\begin{figure}[htb!]
\centering

\caption{Local square relation. A strand with a triple arrow indicates $n-2$ parallel strands oriented according to the direction of the arrow.}
\label{fig:square-lemma}
\end{figure}

\begin{lemma}\label{lem:inductionlargei}
Let $\qroot\in\sC$ such that $q$ is not a $(2k)$-root of unity for any $k=2,3,\dots,n-2$ but allowing $q=\pm1$.  Assuming $\gamma_i\in\basgen$ for all $|i|\leq n-1$, then $\gamma_i\in\basgen$ for all $|i|\geq n$.  
\end{lemma}

\begin{proof}
As a preliminary calculation, let $i\geq n$.  Using the planarizing relation, which is where the restriction on $q$ comes in, write $\gamma_i=cx+dy\in\sk$, where $c=q^{(n-1)/n}$ and $d=-q^{-n(n-1)/2-1/n}/[n-2]!$, and where $x$ and $y$ are the webs appearing in Figure \ref{fig:induction-1of3}. 
\begin{figure}[htb!]
\centering

\caption{The relation $\gamma_i=cx+dy$. A triple arrow indicates $n-2$ parallel strands.}
\label{fig:induction-1of3}
\end{figure}

We see that $x=\gamma_1\gamma_{i-1}\in\sk$.  By another use of the planarizing relation, write $y=cx^\prime+dy^\prime\in\sk$, where $x^\prime$ and $y^\prime$ are the webs appearing in Figure \ref{fig:induction-2of3}.  
 
\begin{figure}[htb!]
\centering

\caption{The relation $y=cx^\prime+dy^\prime$.}
\label{fig:induction-2of3}
\end{figure}

We see that $x^\prime=B_2\gamma_{i-2}\in\sk$ or, in the particular case $n=2$, we have $x^\prime\propto\gamma_{i-2}\in\sk$.  Note that $y^\prime$ is isotopic to the web appearing on the left hand side of Figure \ref{fig:induction-3of3}.  
\begin{figure}[htb!]
\centering

\caption{The relation $y^\prime=s_1\sum_{\sigma\in\sym_{n-2}}b^{\ell(\sigma)}x^{\prime\prime}_\sigma+s_2\sum_{\sigma\in\sym_{n-1}}b^{\ell(\sigma)}y^{\prime\prime}_\sigma$.}
\label{fig:induction-3of3}
\end{figure}

Using Lemma \ref{lem:square-lemma}, write 
\begin{equation*}
y^\prime=s_1\sum_{\sigma\in\sym_{n-2}}b^{\ell(\sigma)}x^{\prime\prime}_\sigma+s_2\sum_{\sigma\in\sym_{n-1}}b^{\ell(\sigma)}y^{\prime\prime}_\sigma\in\sk
\end{equation*}
 where $s_1=a^2a_{n-1}a_{n-2}\prod_{k=2}^{n-1}a_k$ and $s_2=a^2a_{n-1}b\prod_{k=2}^{n-1}a_k$ with $a,b,a_k$ as in Lemma \ref{lem:square-lemma}, and where $x^{\prime\prime}_\sigma$ and $y^{\prime\prime}_\sigma$ are the webs appearing in Figure \ref{fig:induction-3of3}.  We see that $x^{\prime\prime}_\sigma=\rev{\lp\lift{\sigma}_+\rp}\gamma_{i-2}\in\sk$,
where recall the bar notation indicates the reverse orientation.  If $\alpha$ is a cutting arc as in Figure \ref{fig:fig7}, note by Lemma \ref{lem:refinement} that $\rev{\lp\lift{\sigma}_+\rp}\in\sub{\gamma_m;-(n-2)\leq m\leq0}\subset\basgen$, by hypothesis, and $y^{\prime\prime}_\sigma\in\sub{\gamma_m;-(n-2)\leq m\leq i-2}$. 

We now argue by induction on $|i|\geq n$.  For simplicity, assume $i>0$, the $i<0$ case following by a similar argument by reversing orientations and replacing $B_2$ above with $B_{n-2}$.  For the base case, $i=n$, by hypothesis and above we see that: $x=\gamma_1\gamma_{n-1}\in\basgen$; also $x^\prime=B_2\gamma_{n-2}\in\basgen$, or $x^\prime\propto\gamma_{n-2}=\gamma_0\in\basgen$ for $n=2$; also $x^{\prime\prime}_\sigma=\rev{\lp\lift{\sigma}_+\rp}\gamma_{n-2}\in\basgen$; and, $y^{\prime\prime}_\sigma\in\sub{\gamma_m;-(n-2)\leq m\leq n-2}\in\basgen$.  

For the induction step, assume in addition to the hypothesis that $\gamma_j\in\basgen$ for $j=n,n+1,\dots,i-1$. Similar to the base case, the preliminary calculation shows $\gamma_{i}\in\basgen$.  \end{proof}

\end{document}
